\section{Chap 6: Feedback Linearization}
Idea: Does there exist a nonlinear change of variables $z = {[\eta,\xi]}^T$, and a controller input $u = \alpha(x) + \beta(x)v$ that would transform a generic nonlinear SISO system into a partially linear system (\textbf{input-output linearizable}):
\begin{align*}
    \Dot{\eta} & = f_0(\eta,\xi) \\
    \Dot{\xi}  & = A\xi + Bv     \\
    y          & = C\xi
\end{align*}
Or a nonlinear change of variables $z = T(x)$ and a controller input $u = \alpha(x) + \beta(x)v$ resulting in a fully linear system (\textbf{feedback linearizable}):
\begin{align*}
    \Dot{z} & = \Tilde{A}z + \Tilde{B}v
\end{align*}
\subsection{Input Output Linearization}
\textbf{Def}: A function $f$ is called \textbf{smooth} if $f\in C^\infty$ which means $f$ and all of its derivatives are continuous functions.\\
\textbf{Def}: The \textbf{Lie Derivative} $L_f h(x) = \frac{\partial h}{\partial x}f$ along the domain $f$ is defined for SISO systems as:
\begin{align*}
    \Dot{y} = \frac{\partial h}{\partial x} [f(x) + g(x)u] = L_f h(x) + L_g h(x)u
\end{align*}
If the Lie derivative along $g$ is zero $\Dot{y} = L_f h(x) + 0 = L_f h(x)$ which is independent of the input $u$. The process is repeated until the derivative is a function of the input $u$:
\begin{align*}
    y^{(2)} & = \frac{\partial (L_f h(x))}{\partial x}[f(x) + g(x)u]         \\
            & = L_f^2 h(x) + L_g L_f h(x) u                                  \\
    y^{(N)} & = \frac{\partial (L_f^{(N-1)} h(x))}{\partial x}[f(x) + g(x)u] \\
            & = L_f^N h(x) + L_g L_f^{(N-1)} h(x) u
\end{align*}
The input $u$ does only appear in the last equation $y^{N}$ since all previous Lie Derivatives along $g$ were zero. The system is input output linearizable with the feedback controller:
\noindent\begin{equation*}
    \boxed{u = \frac{1}{L_g L_f^{(N-1)}h(x)} \Big(-L_f^N h(x) + v \Big)}
\end{equation*}

This is equal to a chain of $N$ integrators. $N$ is said to be the systems \textbf{relative degree}.

\subsection{Full State Feedback Linearization}
\textbf{Def}: For any two vertor fields $f$ and $g$ on $\mathcal{D} \subset \mathbb{R}^n$ the \textbf{Lie Bracket} $[f,g](x)$ is a third vector field defined as:
\begin{align*}
    [f,g](x) = \frac{\partial g}{\partial x}f(x) - \frac{\partial f}{\partial x}g(x) = L_g f(x) - L_f g(x)
\end{align*}
Taking Lie Brackets of a function is denoted by:
\begin{align*}
    ad_f^0g  & = g(x)                 \\
    ad_f g   & = [f,g](x)             \\
    ad_f^k g & = [f,ad_f^{(k-1)}g](x)
\end{align*}
\textbf{Def}: A distribution $\Delta$ is called \textbf{involutive} iff:
\begin{align*}
    g_1\in \Delta \quad g_2\in \Delta \implies [g_1,g_2] \in \Delta
\end{align*}
\textbf{Theorem}: A SISO system is feedback linearizable iff:
\begin{enumerate}
    \item The matrix $\mathcal{M}(x) = [g(x), ad_f g(x), \dots , ad_f^{(n-1)}g(x)]$ has full column rank ($RANK(\mathcal{M}(x)) = n$)
    \item The distribution $\Delta = SPAN\{g, ad_f g, \dots, ad_f^{(n-2)}g \}$ is involutive
\end{enumerate}
\subsection{Stability and Robustness}
\textbf{Theorem}: The origin is asymptotically stable iff the origin of the zero dynamics $\Dot{\eta} = f(\eta,0)$ are asymptotically stable (minimum phase).\\
The true system will not be known exactly in practice but only approximations $u = \hat{\alpha}(x) + \hat{\beta}(x)K\xi$. This will result in the following system:
\begin{align*}
    \Dot{\eta} & = f(\eta,\xi)                                                                   \\
    \Dot{\xi}  & = (A-BK)\xi + B \delta(z)                                                       \\
    \delta(z)  & = \gamma(x) [(\hat{\alpha}(x)-\alpha(x))                                        \\
               & \quad + (\hat{\beta}(x) - \beta(x) ) KT - \hat{\beta}(x)K(\hat{T}_2(x) - T(x))]
\end{align*}
The uncertain system only differs from the deterministic system by an additive term which means stability can locally be guaranteed.

